\documentclass[11pt]{article}

\usepackage[margin=1in]{geometry}
\usepackage{graphicx}
\usepackage{amsmath}
\usepackage{amssymb}
\usepackage{booktabs}
\usepackage{caption}
\usepackage{float}
\usepackage[hidelinks]{hyperref}

\graphicspath{{figures/}}

\title{Causal Discovery for Materials Degradation:\\[0.35em]
\large A Verified Phase-Field Surrogate and an Intervention-Based Causal Audit}

\author{Fudail ibn Umar Farooq Khawaja$^{1}$ \and Muhammad Ahmad$^{2}$}

\date{%
  $^{1}$Purdue University \qquad
  $^{2}$Thayer School of Engineering, Dartmouth College\\[0.6em]
  Both authors contributed equally.\\[0.6em]
  \today
}

\begin{document}
\maketitle

\begin{abstract}
We built and verified a phase-field solver for the Allen--Cahn equation, trained a Fourier
neural operator (FNO) surrogate for it, and used the resulting synthetic system to study whether
a neural network learning physics relies on the causal drivers of the dynamics or on spurious
correlates. The solver passed a battery of physics-based verification tests: energy monotonicity,
the shrinking-circle curvature law (within $0.8\%$ of theory), and coarsening scaling, and agreed
with two independent numerical schemes. The FNO surrogate reached $1.52\%$ relative $L_2$ error
in distribution and $1.26\%$ on a held-out parameter band. We then constructed a
spurious-correlation dataset in which a nuisance factor is provably non-causal yet correlates with
the outcome in the training environments, and confirmed that a model trained on it adopts the
shortcut, beating a matched control on identical held-out data ($0.038$ vs.\ $0.083$ severity MAE).
A five-diagnostic causal audit, checked against the constructed answer key, cleanly separated
causal from spurious reliance. The nuisance factor carried a Sobol total
sensitivity index of $0.0000$ in the physics but $0.883$ in the model. We report these results
alongside a negative finding from an IRM-based probe and a demonstration that ablating the
identified nuisance direction restores predictive validity on adversarial data.
\end{abstract}

\section{Methods}

\subsection{Allen--Cahn solver and its verification}

We built a phase-field simulator for the Allen--Cahn equation with a spatially varying mobility.
The governing equation is
\begin{equation}
\frac{\partial u}{\partial t} = M(x)\left(\varepsilon^{2}\nabla^{2}u + u - u^{3} + g\right),
\label{eq:allen-cahn}
\end{equation}
where $u$ is the order parameter (the phase field), $M(x)$ a spatially varying mobility,
$\varepsilon$ the interface width, and $g$ a constant tilt of the double-well potential. The
primary solver uses a second-order semi-implicit (SBDF2) pseudo spectral scheme; the variable
mobility is handled with the standard stabilized splitting, treating the constant-coefficient part
implicitly at the maximum mobility $M_{\max}$. We paired the primary solver with two independent
references: an explicit fourth-order Runge--Kutta validator solver and a \texttt{py-pde}
cross-check, together with Nyquist-correct Fourier resampling for changing resolution.

We verified the solver before generating any training data. The verification tests covered energy
monotonicity, the shrinking-circle curvature law $d(R^2)/dt = -2\varepsilon^2$, coarsening scaling,
time-step and grid convergence, and three-way agreement across the two independent schemes; the
outcomes are reported in Section~\ref{sec:solver-results}.

Datasets use Gaussian random fields (GRFs) as the input distribution, chosen because their
correlation length and spectral decay are directly controllable, which yields clean
out-of-distribution holdouts. Figure~\ref{fig:splits} shows one representative example from each
data split: the initial condition $u_0(x)$, the hidden mobility field $M(x)$, and the final state
$u(T)$, with the interface width $\varepsilon$ annotated per split. The main dataset comprises
$6{,}000$ samples at $64\times64$ resolution integrated to $T=1$, together with a full validator
test set re-solved by the independent scheme at $128^2$ resolution and one-quarter the time step.

\begin{figure}[H]
\centering
\includegraphics[width=0.85\textwidth]{fig1_dataset_splits.png}
\caption{One representative sample from each data split (out-of-distribution test, train,
validation, and test). Each row shows the initial condition $u_0(x)$, the hidden mobility field
$M(x)$, and the final state $u(T)$, with the interface width $\varepsilon$ annotated per split.}
\label{fig:splits}
\end{figure}

\subsection{FNO surrogate}

The surrogate is a Fourier neural operator that maps the input fields to the final state. The
architecture was fixed during a heterogeneous-diffusion warm-up on a $3{,}000$-sample dataset,
where $16$ Fourier modes, $64$ hidden channels, and roughly $2.4$M
parameters reached $1.23\%$ in-distribution and $1.20\%$ out-of-distribution relative $L_2$
error, and we reused it unchanged for the Allen--Cahn task. The interface width $\varepsilon$
enters as a constant input channel, following standard parametric-FNO practice, so that the target
is a well-defined function of the inputs. Training used a shared training loop over $100$ epochs;
the loss history is shown in Figure~\ref{fig:training}.

\subsection{Autoencoder featurization}

To provide an analysis substrate for the causal audit, we trained a convolutional autoencoder that
compresses the final-state field into a compact latent representation $z$. The FNO always operates
on full fields; the autoencoder is instrumentation for the audit rather than part of the prediction
pipeline. The encoder uses four stride-$2$ convolution blocks with GroupNorm and GELU activations,
with a mirrored decoder. We swept ten models. Two field variants crossed with latent sizes
$z \in \{8,\dots,128\}$ and evaluated both pixel reconstruction and functional fidelity (the
preservation of physical observables).

\subsection{Constructing the spurious-correlation dataset}

We constructed a dataset with a known ground-truth answer key: a nuisance factor $S$ that is
provably non-causal and a driver $C$ that is provably causal, with training environments in which
$S$ correlates with the outcome and evaluation environments in which the correlation weakens,
breaks, or reverses. The construction hinges on hiding the causal driver: the mobility field $M(x)$
is not observed. If a model saw all causal inputs, the target would be deterministic in them and a
competent network would ignore any nuisance, leaving no shortcut to detect. With $M(x)$ hidden,
the nuisance genuinely carries predictive information in the biased environments, so the
Bayes-optimal predictor must use it, which guarantees both adoption of the shortcut and its
collapse under a reversed correlation. This is the same structure as the ``wolf vs.\ dog''
shortcut-learning problem.

We used the nucleation-growth regime with a tilted well ($g=0.3$) and a hidden mobility of the form
$M(x) = m\cdot\text{GRF}$. Traveling-wave analysis gives the front velocity
\begin{equation}
V = \frac{3\,M\,\varepsilon\,g}{\sqrt{2}},
\label{eq:front-velocity}
\end{equation}
and the critical radius
\begin{equation}
R_c = \frac{\varepsilon\sqrt{2}}{3g},
\label{eq:critical-radius}
\end{equation}
both independent of $m$. Seeding every nucleus above $R_c$ makes severity monotone in the hidden
amplitude for every sample; this held in all twenty of twenty matched-initial-condition
interventions, with severity spanning $0.04$--$0.88$ and no saturation. The full pool contains
$8{,}000$ simulations across seven environments: two biased training environments
($\rho \in \{0.95, 0.8\}$), a control ($\rho = 0$), and four evaluation environments
($\rho \in \{0.95, 0.5, 0, -0.95\}$), accompanied by a per-sample answer-key CSV. From this pool we
trained biased, control, and oracle FNOs, two severity-prediction heads, and the autoencoder later
audited.

\subsection{The causal audit}

We assessed causal versus spurious reliance with five independent diagnostics, checking every
verdict against the constructed answer key:
\begin{enumerate}
\item \emph{Shortcut collapse.} We track the model's error as the correlation between the nuisance
factor and the outcome is progressively shifted and eventually reversed.
\item \emph{Intervention probing.} We intervene separately on the nuisance and causal factors and
measure how strongly the learned latent responds to each.
\item \emph{Invariance analysis (ICP-style).} We test whether the coefficient on each score is
stable across environments; a causal factor's coefficient should be invariant, whereas a spurious
one's should track the correlation strength $\rho$.
\item \emph{Sobol sensitivity.} We compare variance-based total sensitivity indices computed in the
simulator against those computed in the model.
\item \emph{Latent ablation.} We remove the identified nuisance direction from the representation
and measure the effect on adversarial predictive validity and on the decoded physics.
\end{enumerate}

\section{Results}

\subsection{Solver verification}
\label{sec:solver-results}

The verification tests all passed before any training data was generated
(Figure~\ref{fig:solver-verify}). The Ginzburg--Landau energy decays monotonically. The
shrinking-circle test measured a slope of $-1.81\times10^{-3}$ against the theoretical
$-2\varepsilon^2 = -1.80\times10^{-3}$, a deviation of $0.8\%$. The coarsening length obeys
$\ell^2$ linear in $\varepsilon^2 t$, with a fitted slope of $8.9\,\varepsilon^2$ and
$R^2 = 0.979$. Time-step and grid convergence and three-way scheme agreement also held. The choice
of scheme was measured rather than assumed: a first-order semi-implicit stepping missed the
cross-scheme tolerance ($1.8\times10^{-3}$), so we upgraded to the second-order SBDF2 scheme, which
reached $7\times10^{-5}$ at $dt = 2.5\times10^{-3}$.

\begin{figure}[H]
\centering
\includegraphics[width=\textwidth]{fig5_solver_verification.png}
\caption{Solver verification. Left: Ginzburg--Landau energy decay for $M=1$, $g=0$. Middle: the
shrinking-circle test, with the measured slope tracking the theoretical $-2\varepsilon^2$. Right:
coarsening, showing $\ell^2$ growing linearly in $t$ with $R^2 = 0.979$.}
\label{fig:solver-verify}
\end{figure}

\subsection{FNO accuracy}

The surrogate reached $1.52\%$ relative $L_2$ error on the in-distribution test set and $1.53\%$
against targets produced by the independent numerical scheme; because the two schemes agree to
$\sim\!9\times10^{-4}$ on those targets, the surrogate's error budget is essentially all model
rather than solver. On the held-out $\varepsilon$-band it reached $1.26\%$. The per-split errors are
summarized in Figure~\ref{fig:fno-bar}, and a qualitative comparison of solver output, FNO output,
and the pointwise error appears in Figure~\ref{fig:fno-vs-solver}. The training and validation loss
histories are shown in Figure~\ref{fig:training}.

\begin{figure}[H]
\centering
\includegraphics[width=0.75\textwidth]{fig2_fno_error_bar.png}
\caption{Mean relative $L_2$ error of the FNO surrogate across the validation, test, and
independent-scheme (\emph{validator}) target sets, both in-distribution and on the held-out
parameter band (\emph{ood}).}
\label{fig:fno-bar}
\end{figure}

\begin{figure}[H]
\centering
\includegraphics[width=\textwidth]{fig3_fno_vs_solver.png}
\caption{Qualitative FNO accuracy across four samples. Columns show the initial condition $u_0(x)$,
the mobility field $M(x)$, the $\varepsilon$ input channel, the solver final state $u(T)$, the FNO
prediction, and the pointwise error map.}
\label{fig:fno-vs-solver}
\end{figure}

\begin{figure}[H]
\centering
\includegraphics[width=0.75\textwidth]{fig4_training_curve.png}
\caption{Training and validation MSE loss (in normalized target space, log scale) over $100$
training epochs.}
\label{fig:training}
\end{figure}

\subsection{Autoencoder fidelity}

We initially planned a pixel-level reconstruction target, but this is informationally impossible
for saturated two-phase fields: reconstruction error is dominated by interface placement (with
relative $L_2 \approx \sqrt{4\,\delta L}$), so a $5\%$ pixel target would demand roughly
$0.01$-cell interface accuracy from any latent. We therefore re-anchored the fidelity criterion on
what the audit actually requires, preservation of physical observables, while keeping the pixel
curve on the record. At $z = 128$, the reconstructed severity observable reached $R^2 = 0.975$ and
the interface-length observable $R^2 = 0.932$. Representative reconstructions are shown in
Figure~\ref{fig:ae-recon}, and the dependence of observable fidelity on latent dimension in
Figure~\ref{fig:ae-fidelity}.

\begin{figure}[H]
\centering
\includegraphics[width=0.7\textwidth]{fig6_ae_reconstruction.png}
\caption{Autoencoder reconstructions across four samples. Columns show the ground-truth final state
$u(T)$, the reconstruction at latent size $z = 128$, and the error.}
\label{fig:ae-recon}
\end{figure}

\begin{figure}[H]
\centering
\includegraphics[width=0.7\textwidth]{fig7_ae_fidelity.png}
\caption{Functional fidelity of the autoencoder: $R^2$ of the reconstructed severity and
interface-length observables as a function of latent dimension $z$, with reference lines at the
severity and interface targets.}
\label{fig:ae-fidelity}
\end{figure}

\subsection{Adoption of the spurious shortcut}

On held-out same-distribution data, the biased model outperformed its matched control, confirming
that it had taken up the shortcut. The biased FNO reached $0.250$ relative $L_2$ against the
control's $0.351$, and the biased severity head reached $0.038$ MAE against the control's $0.083$,
while the oracle model, which sees the hidden mobility, sat far ahead at $0.0077$. This reproduces
the full designed ordering of oracle, biased, and control models.

\subsection{Causal-audit findings}

All five diagnostics agreed with the constructed answer key, and their principal results are
collected in Table~\ref{tab:diagnostics}. The single most telling number is the variance-based
sensitivity of the nuisance factor: its Sobol total index is $0.0000$ in the physics but $0.883$ in
the model, while the hidden causal driver carries an index of $0.644$ in the physics and $0.0000$
in the model. In other words, the model's sensitivity structure is the mirror image of the true
physics.

\begin{table}[H]
\centering
\caption{Principal result of each of the five audit diagnostics.}
\label{tab:diagnostics}
\begin{tabular}{@{}p{0.28\textwidth}p{0.64\textwidth}@{}}
\toprule
\textbf{Diagnostic} & \textbf{Key result} \\
\midrule
Shortcut collapse &
Biased head error rises $327\%$ under a flipped correlation ($0.038 \to 0.161$, monotone in the
$\rho$-shift); the control changes by $-4\%$ and the oracle stays flat at $0.0073$. \\
\addlinespace
Intervention probing &
The latent responds to nuisance interventions with $R^2 = 0.976$ and to causal interventions with
$R^2 = 0.495$; the two response directions sit $80.4^\circ$ apart. \\
\addlinespace
Invariance (ICP-style) &
The causal-score coefficient is stable across environments (relative range $0.145$), whereas the
nuisance-score coefficient tracks $\rho$ (a $\sim\!222\times$ swing, sign-flipping on the flipped
data). \\
\addlinespace
Sobol sensitivity &
Nuisance total index $0.0000$ (simulator) vs.\ $0.883$ (model); hidden driver $0.644$ (simulator)
vs.\ $0.0000$ (model). \\
\addlinespace
Latent ablation &
Removing the nuisance direction moves the evaluation-ID probe $R^2$ from $0.868$ to $0.292$ and the
flipped-environment probe $R^2$ from $-1.028$ to $+0.470$, while the decoded physics channel
changes by only $4.3\%$. \\
\bottomrule
\end{tabular}
\end{table}

Beyond detection, the latent-ablation diagnostic doubles as a mitigation: surgically removing the
identified nuisance direction from the learned representation restored predictive validity on
adversarial data, with the probe $R^2$ moving from $-1.03$ to $+0.47$ while leaving the physical
content roughly $96\%$ intact.

One secondary check produced a negative result. An IRMv1-penalized probe generalized \emph{worse}
than pooled empirical risk minimization (ERM) at every penalty weight, reaching $0.289$ against
ERM's $0.425$ flipped-environment $R^2$, with the penalty weight $\lambda$ selected on an
independent shifted environment. Two factors plausibly account for this. The pooled baseline is
already implicitly invariance-regularized, since half of its data comes from the $\rho = 0$
environment, and linear IRMv1 is documented to underperform in settings with low environment
diversity~\cite{rosenfeld2021risks}. We report the result as measured rather than tuning it away.

\section{Discussion}

\subsection{Technical challenges and their resolution}

Several difficulties shaped the final design. The first surrogate, trained on a $T=4$ horizon,
memorized its training set (train MSE $5.5\times10^{-4}$) yet generalized at only $17\%$, with its
best validation loss occurring at epoch $6$ of $100$: a $T=4$ one-shot map spans an entire
coarsening cascade whose merger outcomes depend sensitively on the input, so the operator is poorly
conditioned. We shortened the horizon to $T=1$, where the operator is well-conditioned, and left
longer horizons via autoregressive composition to future work.

Even at $T=1$ the surrogate showed $11\%$ error with immediate overfitting. The root cause was that
the parameter $\varepsilon$ was not a model input and is not inferable from the other channels, so
the target was not, strictly, a function of the inputs. Feeding $\varepsilon$ as a constant channel
collapsed the error roughly sixtyfold to $1.5\%$. The same hidden-variable structure that was a bug
here is precisely what we later build on purpose in the spurious-correlation dataset, where its
early-plateau training signature reappears by design.

The physics checks initially targeted the textbook coarsening exponent $\ell \sim t^{1/2}$, which is
unmeasurable in a unit box: domains reach a quarter of the box within $t \approx 1$, leaving no
scaling decade, and the structure-factor length is further biased by the Porod tail. We restated the
check in the reachable regime as the differential form of the same law, $\ell^2$ linear in
$\varepsilon^2 t$, which we measured at $R^2 = 0.98$ with slope $\approx 9\varepsilon^2$ using an
interface-density length measure.

Spectral resampling also required care. Naive Fourier truncation and zero-padding broke the exact
up/down-sample round trip (a $2.6\%$ error) by mishandling the ambiguous Nyquist bin. We replaced it
with Nyquist-splitting resampling matrices, in which the coefficient is split on upsampling and
re-summed on downsampling, after which the round trip is exact and the validator targets are
alias-free.

Finally, compute was a practical constraint. A single-threaded NumPy implementation projected about
$14$ hours for one dataset plus its validator set. Moving both solvers to float32 CUDA backends gave
roughly a $100\times$ speedup, cutting the total to about $25$ minutes, while the float64 NumPy path
was retained as the verified reference with cross-backend agreement unit-tested to below
$10^{-4}$. This also made the audit's $1{,}280$ additional Sobol simulations inexpensive. Along the
way we worked around one ecosystem breakage: SALib 1.5.1 passes its problem names into
\texttt{pd.unique}, which pandas 3.0 no longer accepts for plain lists, fixed by passing the names
as an ndarray.

Two audit checks were initially mis-formalized, and we corrected both without touching any
threshold. First, requiring the multivariate probe's causal coefficient to be stable across
environments is not what invariant causal prediction predicts: with a shortcut available, a
full-representation fit will legitimately shift weight from the causal to the spurious direction,
and that reweighting is itself evidence of adoption. We replaced this with ICP-correct univariate
score probes, which pass cleanly, and kept the multivariate instability as corroborating evidence.
Second, the IRM check originally tested a coefficient ratio rather than IRM's operative promise,
out-of-environment risk, so we re-formalized it with $\lambda$ selected on an independent
shifted environment, and it remained negative, which is the finding reported above.

\subsection{Scope of the claims}

Two claims are supported by the evidence above: that the surrogate is accurate enough to stand in
for the solver, and that the audit separates causal from spurious reliance. A third claim, that this
PDE describes real materials degradation, is deliberately not made. Establishing it requires
experimental data, which is the subject of ongoing work.

\subsection{Limitations}

All headline metrics come from a single master seed per phase; a multi-seed replication with
dispersion estimates should precede any publication claim. The surrogate covers $T = 1$ only, so
autoregressive composition to longer horizons remains unvalidated. The mitigation demonstrated here
is post-hoc, operating on a trained representation; training-time mitigation was probed only through
the inconclusive IRM linear probe. Finally, the audit has so far been run on a single governing
equation, so replication on a dynamically distinct system, such as Fisher--KPP, is the most
important outstanding test of generality.

\section*{Acknowledgements}

This work was conducted at CEDAR Lab, Thayer School of Engineering, Dartmouth College, advised by
Bijan Mazaheri.

\begin{thebibliography}{1}
\bibitem{rosenfeld2021risks}
E.~Rosenfeld, P.~Ravikumar, and A.~Risteski.
\newblock The risks of invariant risk minimization.
\newblock In \emph{International Conference on Learning Representations (ICLR)}, 2021.
\end{thebibliography}

\end{document}
